% 01-architecture.tex — Chapter 1: Architecture and Pipeline
\chapter{Architecture and Pipeline}
\label{chap:01-architecture}
\index{architecture}
\index{pipeline}

\pnum
The Lain compiler is a single-pass, single-binary compiler written entirely in C99.
Every subsystem is implemented as a header file (\texttt{.h}); the single translation
unit \srcref{src/main.c} includes them all in dependency order.  This design eliminates
link-time symbol management, allows the C compiler to inline aggressively across
subsystem boundaries, and makes the build reproducible with a single invocation of
\texttt{gcc} or \texttt{clang}.

\pnum
The compiler targets C99 as its sole output language.  Lain source text is translated
to a single C file that is compiled by the host C toolchain.  This keeps the Lain
compiler itself small and portable: it relies on the host compiler for platform ABI
handling, optimisation, and linking.

% -------------------------------------------------------------------------
\section{Design philosophy}
\label{sec:arch-philosophy}
\index{design philosophy}
% -------------------------------------------------------------------------

\pnum
Four principles shape every architectural decision in the compiler.

\begin{description}[leftmargin=2cm, style=nextline]
\item[Single translation unit.]
  All \texttt{.h} files are included into \srcref{src/main.c}.  There is no
  separate compilation, no static library, and no linker step for the compiler
  itself.  The include order is a direct expression of the dependency graph.
\item[Three arenas, zero free.]
  Memory is allocated from three bump-pointer arenas (see \cref{sec:arch-arenas}).
  No \texttt{malloc}/\texttt{free} calls appear in compiler logic; all memory is
  released in bulk when the compiler exits.  This eliminates entire classes of
  memory-safety bugs and makes the compiler trivially valgrind-clean.
\item[Two semantic phases, one pass each.]
  Semantic analysis runs in exactly two ordered phases per function: (1)~name
  resolution and type inference (\textit{resolve}), and (2)~ownership, bounds,
  and purity checks (\textit{walk}).  The flag \cname{sema\_walk\_phase} separates
  them.  No phase reads results produced by a later phase.
\item[C99 backend.]
  Emitting C99 instead of machine code delegates register allocation, instruction
  selection, and ABI handling to the host toolchain.  The tradeoff is accepted
  compilation latency in exchange for a dramatically smaller compiler.
\end{description}

% -------------------------------------------------------------------------
\section{Compilation pipeline}
\label{sec:arch-pipeline}
\index{pipeline!compilation stages}
% -------------------------------------------------------------------------

\pnum
\Cref{fig:pipeline} shows the compiler pipeline.  Each stage is a self-contained
module; data flows forward through the pipeline.

\begin{figure}[ht]
\centering
\begin{tikzpicture}[node distance=0.9cm and 1.6cm]
  % Source
  \node[databox] (src) {source.ln};

  % load_module
  \node[passbox, right=of src] (load) {load\_module};

  % DeclList
  \node[databox, right=of load] (ast) {DeclList (AST)};

  % sema
  \node[passbox, below=of ast] (sema) {sema\_resolve\\\_module};

  % emit
  \node[passbox, left=of sema] (emit) {emit};

  % out.c
  \node[databox, left=of emit] (outc) {out.c};

  % gcc
  \node[passbox, below=of outc] (gcc) {host gcc/clang};

  % binary
  \node[databox, right=of gcc] (bin) {binary};

  % Arrows top row
  \draw[dataarrow] (src)  -- (load);
  \draw[dataarrow] (load) -- (ast);
  \draw[dataarrow] (ast)  -- (sema);

  % Arrows middle row
  \draw[dataarrow] (sema) -- (emit);
  \draw[dataarrow] (emit) -- (outc);

  % Arrows bottom row
  \draw[dataarrow] (outc) -- (gcc);
  \draw[dataarrow] (gcc)  -- (bin);

  % Labels
  \node[font=\tiny\itshape, below=0.1cm of load] {lex + parse + import};
  \node[font=\tiny\itshape, below=0.1cm of sema] {resolve, infer, check};
  \node[font=\tiny\itshape, below=0.1cm of emit] {C99 codegen};
\end{tikzpicture}
\caption{Lain compilation pipeline.}
\label{fig:pipeline}
\end{figure}

\pnum
The five major stages are:

\begin{enumerate}
\item \textbf{Module loading} (\srcref{src/module.h}).  Given the root module name,
  \cname{load\_module()} reads the source file from disk using \cname{file\_arena},
  lexes it with \cname{lexer\_next()}, parses it with the recursive-descent parser,
  and recursively loads any \texttt{import}ed modules.  The result is a singly-linked
  \cname{DeclList} of top-level declarations allocated in \cname{ast\_arena}.

\item \textbf{Semantic analysis} (\srcref{src/sema.h} and \srcref{src/sema/}).
  \cname{sema\_resolve\_module()} runs two ordered phases on every function body:
  \begin{enumerate}[(a)]
  \item \textit{Resolve phase}: \cname{sema\_resolve\_stmt()} populates identifier
    \cname{->decl} pointers and infers expression types.  UFCS desugaring,
    comptime generic substitution, and destructuring expansion happen here.
  \item \textit{Walk phase}: \cname{walk\_stmt()} runs VRA, bounds checking,
    ownership/linearity checking, purity checking, and bounded-while verification.
    The flag \cname{sema\_walk\_phase} is \texttt{true} during this phase only.
  \end{enumerate}
  After both phases complete for a function, \cname{sema\_check\_function\_linearity()}
  verifies that all owned values have been consumed or moved.

\item \textbf{Mutual recursion check} (\srcref{src/sema.h:\,786}).
  After all functions are resolved, \cname{sema\_check\_no\_mutual\_recursion()}
  performs DFS over the call graph and rejects any cycle that involves only
  \texttt{func} (pure) declarations.

\item \textbf{Code generation} (\srcref{src/emit.h} and \srcref{src/emit/}).
  \cname{emit()} traverses the annotated \cname{DeclList} and writes valid C99 to
  the output file.  Struct types are forward-declared in topological order;
  constructors are synthesised for enum variants; slice types receive generated
  \texttt{typedef} names.

\item \textbf{Host compilation.}  The compiler invokes no subprocess; it is the
  user's responsibility to compile the generated C file.  The recommended invocation
  is documented in \cref{sec:arch-build}.
\end{enumerate}

% -------------------------------------------------------------------------
\section{Arena memory model}
\label{sec:arch-arenas}
\index{arena allocator}
\index{memory!arenas}
% -------------------------------------------------------------------------

\pnum
The compiler uses three bump-pointer arenas declared in \srcref{src/main.c}:

\begin{center}
\begin{tabular}{lll}
\toprule
\textbf{Arena variable} & \textbf{Contents} & \textbf{Lifetime} \\
\midrule
\cname{file\_arena}   & Source text buffers, raw file bytes & Entire compiler run \\
\cname{ast\_arena}    & All AST nodes (\cname{Expr}, \cname{Stmt}, \cname{Decl},
                         \cname{Type}), & Entire compiler run \\
                      & identifier strings, module records & \\
\cname{\_sema\_arena} & Symbol tables, \cname{RangeTable}, borrow entries, & Entire compiler run \\
                      & in-guard entries, constraint lists & \\
\bottomrule
\end{tabular}
\end{center}

\begin{invariant}
No \cname{free()} call is made on arena-allocated memory during compilation.
Arenas are released only when the process exits.  Compiler passes must never
store a raw pointer to stack-allocated data inside arena-allocated structures.
\end{invariant}

\pnum
The arena implementation (\srcref{src/utils/arena.h}) is a three-field struct
\cname{\{beg, cur, end\}}.  Allocation advances \cname{cur} by the requested
size with optional alignment padding.  When \cname{ARENA\_DEBUG} is enabled (the
default), every push asserts that sufficient space remains, making out-of-memory
conditions loud rather than silent.

\pnum
The allocation macros are:

\begin{center}
\begin{tabular}{ll}
\toprule
\textbf{Macro} & \textbf{Usage} \\
\midrule
\cname{arena\_push(a, T)}               & Allocate one \texttt{T} \\
\cname{arena\_push\_many(a, T, n)}       & Allocate \texttt{n} consecutive \texttt{T} values \\
\cname{arena\_push\_aligned(a, T)}       & Allocate one \texttt{T} with \texttt{alignof(T)} padding \\
\cname{arena\_push\_many\_aligned(a,T,n)}& Aligned array of \texttt{n} \texttt{T} values \\
\bottomrule
\end{tabular}
\end{center}

% -------------------------------------------------------------------------
\section{Source file structure}
\label{sec:arch-files}
\index{source files}
% -------------------------------------------------------------------------

\pnum
\Cref{tab:srcfiles} lists every header file in the compiler source tree together
with its approximate line count and role.

\begin{longtable}{lrl}
\caption{Compiler source files.\label{tab:srcfiles}} \\
\toprule
\textbf{File} & \textbf{LOC} & \textbf{Role} \\
\midrule
\endfirsthead
\multicolumn{3}{l}{\small (continued)} \\
\toprule
\textbf{File} & \textbf{LOC} & \textbf{Role} \\
\midrule
\endhead
\bottomrule
\endfoot
\srcref{src/main.c}                  &  115 & Entry point, arena init, pipeline orchestration \\
\srcref{src/token.h}                 &  340 & \cname{TokenKind} enum, \cname{Token} struct \\
\srcref{src/lexer.h}                 &  384 & FSM lexer, \cname{LexerState}, \cname{lexer\_next()} \\
\srcref{src/ast.h}                   & 1032 & All AST node types (\cname{Expr}, \cname{Stmt}, \cname{Decl}, \cname{Type}) \\
\srcref{src/ast\_clone.h}            &  327 & Deep-clone of AST subtrees (used by generics) \\
\srcref{src/ast\_print.h}            &  508 & \texttt{-{}-dump-ast} debug printer \\
\srcref{src/parser.h}                &   -- & Top-level parser include \\
\srcref{src/parser/core.h}           &  199 & Parser state, lookahead, \cname{expect()}, error recovery \\
\srcref{src/parser/expr.h}           &  403 & Expression parser (Pratt precedence climbing) \\
\srcref{src/parser/stmt.h}           &  736 & Statement parser \\
\srcref{src/parser/decl.h}           &  801 & Declaration parser (func, proc, struct, enum, \ldots) \\
\srcref{src/parser/type.h}           &   -- & Type expression parser \\
\srcref{src/module.h}                &  142 & Module loading, import resolution, \cname{ModuleNode} \\
\srcref{src/sema.h}                  & 1066 & Sema orchestrator, VRA, bounded-while verifier \\
\srcref{src/sema/scope.h}            &  146 & Lexical scope (global/local symbol tables) \\
\srcref{src/sema/resolve.h}          &  987 & Phase 1: name resolution, UFCS desugaring \\
\srcref{src/sema/typecheck.h}        & 1062 & Phase 2 type inference, integer widening \\
\srcref{src/sema/bounds.h}           &   -- & Bounds-check emission helpers \\
\srcref{src/sema/ranges.h}           &  594 & VRA: \cname{Range}, \cname{RangeTable}, interval arithmetic \\
\srcref{src/sema/region.h}           &  538 & Borrow checker: two-phase borrows, \cname{BorrowTable} \\
\srcref{src/sema/linearity.h}        & 1398 & Ownership/linearity checker, \cname{LinearEnv} \\
\srcref{src/sema/generic.h}          &  216 & Comptime generic instantiation \\
\srcref{src/sema/comptime.h}         &  321 & Comptime expression evaluator \\
\srcref{src/sema/exhaustiveness.h}   &  315 & Pattern exhaustiveness checker \\
\srcref{src/sema/use\_analysis.h}    &  260 & Use-before-init / unused-variable analysis \\
\srcref{src/emit.h}                  &   -- & Top-level emit include \\
\srcref{src/emit/core.h}             &  370 & Emit state, \cname{EmitCtx}, file writer helpers \\
\srcref{src/emit/type\_order.h}      &  203 & Topological sort of struct types \\
\srcref{src/emit/decl.h}             &  523 & Declaration emitter (functions, structs, enums) \\
\srcref{src/emit/stmt.h}             &  627 & Statement emitter \\
\srcref{src/emit/expr.h}             &  700 & Expression emitter \\
\srcref{src/emit/ctor.h}             &   -- & Enum constructor function emitter \\
\srcref{src/emit/lain\_header.h}     &  325 & Lain runtime header injected at top of output \\
\srcref{src/utils/arena.h}           &  124 & Bump-pointer arena allocator \\
\srcref{src/utils/common.h}          &   -- & Common primitive types (\cname{isize}, \cname{usize}, \ldots) \\
\srcref{src/utils/ascii.h}           &  241 & ASCII classification helpers \\
\srcref{src/utils/file.h}            &   -- & Platform file I/O (\cname{file\_read\_all}) \\
\srcref{src/utils/panic.h}           &   -- & \cname{panic\_if\_msg()} assertion macro \\
\srcref{src/utils/string.h}          &   -- & String utilities \\
\srcref{src/utils/utf8.h}            &   -- & UTF-8 helpers \\
\srcref{src/args.h}                  &   -- & Command-line argument parsing \\
\end{longtable}

% -------------------------------------------------------------------------
\section{Naming conventions}
\label{sec:arch-naming}
\index{naming conventions}
% -------------------------------------------------------------------------

\pnum
Function and variable names in the compiler follow consistent prefixes that
identify the owning subsystem.

\begin{center}
\begin{tabular}{ll}
\toprule
\textbf{Prefix} & \textbf{Subsystem} \\
\midrule
\cname{lexer\_}     & Lexer (\srcref{src/lexer.h}) \\
\cname{parse\_}     & Parser (\srcref{src/parser/}) \\
\cname{sema\_}      & Semantic analysis orchestrator (\srcref{src/sema.h}) \\
\cname{sema\_resolve\_} / \cname{resolve\_} & Name resolution (\srcref{src/sema/resolve.h}) \\
\cname{sema\_infer\_}  & Type inference (\srcref{src/sema/typecheck.h}) \\
\cname{range\_}        & Value range analysis (\srcref{src/sema/ranges.h}) \\
\cname{borrow\_}       & Borrow checker (\srcref{src/sema/region.h}) \\
\cname{linear\_} / \cname{sema\_check\_} & Linearity checker (\srcref{src/sema/linearity.h}) \\
\cname{emit\_}         & Code generator (\srcref{src/emit/}) \\
\cname{arena\_}        & Arena allocator (\srcref{src/utils/arena.h}) \\
\cname{mrec\_}         & Mutual-recursion DFS (\srcref{src/sema.h}) \\
\cname{comptime\_}     & Comptime evaluator (\srcref{src/sema/comptime.h}) \\
\bottomrule
\end{tabular}
\end{center}

\pnum
Type names follow the pattern \textit{SubsystemConcept}, e.g.\ \cname{RangeTable},
\cname{BorrowEntry}, \cname{LinearEnv}, \cname{InGuardEntry}.  List nodes use the
suffix \cname{List} (e.g.\ \cname{DeclList}, \cname{StmtList}, \cname{ExprList}).

% -------------------------------------------------------------------------
\section{Global state}
\label{sec:arch-globals}
\index{global state}
% -------------------------------------------------------------------------

\pnum
The compiler maintains a small set of global variables that constitute the shared
state visible across subsystem headers.  These are declared (not merely declared
\texttt{extern}) in the header that owns them; because all headers are included
into a single translation unit, each global is defined exactly once.

\begin{longtable}{lll}
\caption{Compiler global state.\label{tab:globals}} \\
\toprule
\textbf{Variable} & \textbf{Type} & \textbf{Purpose} \\
\midrule
\endfirsthead
\multicolumn{3}{l}{\small (continued)} \\
\toprule
\textbf{Variable} & \textbf{Type} & \textbf{Purpose} \\
\midrule
\endhead
\bottomrule
\endfoot
\cname{sema\_source\_text}   & \cname{const char *} & Pointer into source buffer; used by \cname{diagnostic\_show\_line} \\
\cname{sema\_source\_file}   & \cname{const char *} & Source file name for error messages \\
\cname{sema\_in\_guards}     & \cname{InGuardEntry *} & Head of \texttt{in}-guard linked list \\
\cname{sema\_decls}          & \cname{DeclList *}   & Top-level declarations of current module \\
\cname{sema\_arena}          & \cname{Arena *}      & Arena for sema-phase allocations \\
\cname{sema\_ranges}         & \cname{RangeTable *} & VRA range table for current function \\
\cname{sema\_in\_unsafe\_block} & \cname{bool}      & \texttt{true} inside an \texttt{unsafe \{...\}} block \\
\cname{sema\_walk\_phase}    & \cname{bool}         & \texttt{true} during walk phase only \\
\cname{current\_return\_type}  & \cname{Type *}     & Return type of function being analysed \\
\cname{current\_function\_decl} & \cname{Decl *}   & \cname{Decl} of function being analysed \\
\cname{current\_module\_path} & \cname{const char *}& Path of module being analysed \\
\cname{emit\_source\_filename} & \cname{const char *}& Source filename embedded in emit output \\
\bottomrule
\end{longtable}

\begin{note}
The globals listed above are declared in \srcref{src/sema.h} and
\srcref{src/emit/core.h}.  Because they are file-scope statics or non-static
definitions inside a single translation unit, no header guard is needed for
uniqueness.  The include order enforced by \srcref{src/main.c} ensures that
consumer headers can reference them.
\end{note}

% -------------------------------------------------------------------------
\section{Build instructions}
\label{sec:arch-build}
\index{build}
% -------------------------------------------------------------------------

\pnum
The compiler is built with a single compilation command.  No build system is
required.

\begin{ccode}
gcc -std=c99 -Wall -Wextra -o lain src/main.c -I src
\end{ccode}

\pnum
The generated C output is compiled by the host toolchain.  The runtime header
injected at the top of every generated file uses \texttt{\#define} aliases for
standard C library functions (e.g.\ \texttt{libc\_printf}).  When compiling
generated output directly, pass the appropriate \texttt{-D} flags:

\begin{ccode}
gcc -o out out.c -Dlibc_printf=printf -Dlibc_puts=puts -w
\end{ccode}

\pnum
The \texttt{-{}-dump-ast} flag instructs the compiler to print the parsed AST to
\texttt{stdout} and exit without running sema or emit.  This is useful for
debugging parser changes.

\begin{ccode}
./lain --dump-ast source.ln
\end{ccode}
